\documentclass[11pt,a4paper]{article}

\usepackage[utf8]{inputenc}
\usepackage[T1]{fontenc}
\usepackage[spanish,es-nodecimaldot,es-noshorthands]{babel}
\usepackage{lmodern}
\usepackage[margin=2.5cm]{geometry}
\usepackage{amsmath,amssymb,amsthm}
\usepackage{booktabs,array,longtable}
\usepackage{enumitem}
\usepackage{xcolor}
\usepackage[colorlinks=true,linkcolor=blue!50!black,citecolor=blue!50!black,urlcolor=blue!50!black]{hyperref}

% ---------------------------------------------------------------
% Entornos
% ---------------------------------------------------------------
\theoremstyle{plain}
\newtheorem{proposicion}{Proposición}
\newtheorem{conjetura}{Conjetura}
\newtheorem{corolario}{Corolario}
\theoremstyle{definition}
\newtheorem{definicion}{Definición}
\newtheorem{hipotesis}{Hipótesis}
\theoremstyle{remark}
\newtheorem{observacion}{Observación}

\newcommand{\Nocc}{N_{\mathrm{occ}}}
\newcommand{\vac}{\varnothing}
\newcommand{\R}{\mathbb{R}}
\newcommand{\one}{\mathbf{1}}
\newcommand{\Hcl}{H^{\mathrm{cl}}}
\newcommand{\Icl}{I^{\mathrm{cl}}}
\newcommand{\modP}{\textsf{P}}
\newcommand{\modS}{\textsf{S}}
\newcommand{\pendiente}[1]{\textcolor{red!70!black}{[#1]}}

\newenvironment{fase}[1]{\subsection*{#1}\begin{description}[leftmargin=0pt,itemsep=2pt]}{\end{description}}

\title{Una extensión matricial del modelo de Schelling:\\
lectura crítica del borrador y programa de investigación}
\author{Documento de trabajo del proyecto \emph{Modelos Basados en Agentes}}
\date{29 de septiembre de 2026}

\begin{document}
\maketitle

\begin{abstract}
Este documento tiene dos partes. En la primera se hace una lectura crítica del
borrador \emph{Una extensión matricial del modelo de Schelling basada en la
incomodidad local} (versión~8) y del código Mesa asociado. La conclusión
principal es que el borrador desarrolla con cuidado una sola de las dos
versiones del modelo que nos interesan (la \emph{promediada}), mientras que el
código implementa la otra (la \emph{de suma}), y que en el caso de dos
poblaciones la versión promediada es exactamente equivalente a un modelo de
Schelling clásico con orientación y umbral propios de cada población: la
riqueza genuina del enfoque matricial aparece en la versión de suma, en el
tratamiento de las casillas vacías y en el caso de tres o más poblaciones. En la
segunda parte se propone un marco unificado para ambas versiones y un programa
de investigación en nueve fases, con hipótesis contrastables, observables,
protocolo de simulación y criterios de cierre para cada fase.
\end{abstract}

\tableofcontents

% ===============================================================
\section{Propósito y alcance}
% ===============================================================

El objetivo del proyecto es describir \emph{cuantitativamente} la dinámica de
dos generalizaciones del modelo de desplazamiento de Schelling, gobernadas por
una matriz de coeficientes de afinidad--aversión $H$ y un umbral de tolerancia
$\tau$:
\begin{itemize}
  \item \textbf{Modelo \modP{} (promediado).} La incomodidad del agente es el
  promedio de los coeficientes $H_{t_a,t_b}$ sobre sus vecinos ocupados; es
  decir, depende de las \emph{proporciones} de cada población en la vecindad.
  Es el modelo del borrador.
  \item \textbf{Modelo \modS{} (de suma).} La incomodidad es la suma de los
  coeficientes sobre los vecinos ocupados; depende de los \emph{recuentos}.
  Es el modelo implementado en \texttt{schelling\_general2.py},
  \texttt{schelling\_general3.py}, \texttt{schelling\_red2.py} y, como caso
  particular, en \texttt{modelo.py} (modelo de tres grupos con población
  neutral).
\end{itemize}
Ambas dinámicas dependen además de la geometría: retículos toroidales con
vecindades de Moore o de von Neumann de radio $r$, y eventualmente redes.

El documento está pensado para usarse como hoja de ruta: cada fase de la
sección~\ref{sec:programa} tiene un objetivo, tareas, un entregable y un
criterio de cierre, de modo que podamos avanzar paso a paso y revisar el plan al
final de cada fase.

% ===============================================================
\section{Lectura crítica del borrador}
\label{sec:critica}
% ===============================================================

\subsection{Lo que el borrador establece bien}

\begin{enumerate}[label=(\alph*)]
  \item La definición de la incomodidad local inducida por $H$ en la versión
  promediada y la recuperación del modelo clásico con
  $\Hcl=\left(\begin{smallmatrix}0&1\\1&0\end{smallmatrix}\right)$.
  \item La invariancia de la regla de mejora bajo transformaciones afines
  positivas $\widetilde H=\alpha H+\beta\one$, $\alpha>0$, con el umbral
  transformado $\widetilde\tau=\alpha\tau+\beta$ (ecuación~(1) del borrador).
  \item La observación de que el efecto de un movimiento sobre la incomodidad
  total es local, y de que la mejora individual no implica, en general,
  disminución de la incomodidad total.
  \item La descomposición por filas del caso general $2\times2$
  (sección~5.1), que es la observación clave: la orientación de la población
  $i$ depende solo del signo del contraste $c_i=H_{i,\bar\imath}-H_{i,i}$.
  \item La intuición de que las fronteras $c_i=0$ separan regímenes
  cualitativamente distintos y de que conviene estudiar observables
  macroscópicos (segregación, agentes cómodos, tiempo de convergencia,
  clústeres).
\end{enumerate}

\subsection{Observaciones de fondo}

\paragraph{O1. Dos modelos, una sola teoría.}
El borrador desarrolla únicamente la versión promediada, pero todo el código del
proyecto implementa la versión de suma. Las dos no son equivalentes salvo a
densidad máxima (sin vacíos), y sus propiedades matemáticas difieren en puntos
esenciales (grupo de invariancia, papel de los vacíos, existencia de función
potencial; véase la sección~\ref{sec:marco}). Además, no hay en el repositorio
una implementación de la versión promediada, por lo que no está claro con qué
programa se generaron las figuras~1--11. El programa debe tratar ambos modelos
en paralelo y con una única implementación parametrizada.

\paragraph{O2. En el caso de dos poblaciones, el modelo \modP{} se reduce al
clásico.}
La sección~5.1 del borrador contiene ya el resultado, pero no se extrae la
consecuencia. Para un agente de tipo $i$ con $n_a>0$,
\[
  I^H_a = H_{i,i} + c_i\,\Icl_a ,\qquad c_i = H_{i,\bar\imath}-H_{i,i},
\]
y la condición de comodidad con el umbral común $\tau$ es
$\Icl_a\le\theta_i:=(\tau-H_{i,i})/c_i$ si $c_i>0$ (y la desigualdad opuesta si
$c_i<0$). Por tanto la dinámica completa queda determinada, para cada población,
por una orientación (homofílica/heterofílica) y un umbral clásico efectivo
$\theta_i$ (Proposición~\ref{prop:reduccion}). Consecuencias:
\begin{itemize}
  \item El modelo \modP{} de dos poblaciones con matriz general y $\tau$ común
  \emph{no es más general} que un modelo de Schelling clásico con umbrales
  distintos por población y con la posibilidad de invertir la preferencia. Esto
  lo conecta directamente con la literatura sobre poblaciones heterogéneas
  \cite{HatnaBenenson2015,AcunaGoles2025}.
  \item La normalización $(\varepsilon,\sigma,\rho)$ de las secciones~4.2 y
  5.2 mezcla filas y no es la parametrización natural; la parametrización
  natural es $(s_i,\theta_i)_{i}$, con $s_i=\operatorname{sgn}c_i$.
  \item La asimetría intergrupal $A=\beta-\delta$ (sección~5.3) \emph{no tiene
  contenido dinámico propio} en \modP: dos matrices con los mismos
  $(s_i,\theta_i)$ producen exactamente la misma dinámica aunque tengan
  asimetrías distintas. Sí tiene contenido en \modS{} (O-S más abajo).
  \item Las secciones 4 y 5 pueden condensarse en una proposición y una tabla,
  liberando espacio para lo que es nuevo.
\end{itemize}

\paragraph{O3. La convención del vecindario vacío rompe la invariancia afín.}
Con la convención $I^H_a=0$ cuando $n_a=0$, una casilla aislada vale $0$ bajo
$H$ pero también $0$ bajo $\widetilde H=\alpha H+\beta\one$, cuando debería
valer $\beta$ para que la transformación sea una invariancia. Así, en cuanto hay
casillas aisladas disponibles como destino, $H$ y $\widetilde H$ \emph{no}
inducen la misma dinámica. El borrador lo detecta para la incomodidad total
(observación sobre vecindarios vacíos) pero no para la regla de movimiento, que
es lo relevante. La propia duda del texto (``no sé si es mejor tomar
$\Icl_a=1$'') muestra que la convención es en realidad un \emph{parámetro} del
modelo. Propuesta: introducir el valor de aislamiento $\iota_i$ de cada tipo, o
bien, de forma más elegante, añadir a $H$ una columna para el vacío $H_{i,\vac}$
(sección~\ref{sec:marco}). Esto además unifica \modP{} y \modS.

\paragraph{O4. Las rectas $c_i=0$ no son bifurcaciones en sentido estricto.}
Al acercarse a $c_i=0$ con $\tau$ fijo, $\theta_i\to\pm\infty$: la población
queda \emph{congelada} (todos cómodos) si $\tau>H_{i,i}$, o en \emph{búsqueda
perpetua} (todos incómodos, sin mejora posible salvo el aislamiento) si
$\tau<H_{i,i}$. Es decir, antes de llegar a la recta neutral el sistema atraviesa
regímenes triviales determinados por el umbral. Los cambios macroscópicos
interesantes ocurren a valores finitos de $\theta_i\in(0,1)$, que es donde hay
que buscar transiciones (como en el diagrama de fases de
\cite{Gauvin2009}). Esto explica, por ejemplo, la figura~3 y la lentitud
observada en el caso $\alpha=0{,}4$, $\beta=\gamma=0{,}6$, $\tau=0{,}5$: allí la
población $-1$ es neutral con incomodidad constante $0{,}6>\tau$, está siempre
incómoda y solo puede mejorar yendo a una casilla aislada, de modo que en cada
paso cada agente $-1$ recorre toda la retícula buscando una.

\paragraph{O5. Los umbrales de las figuras caen exactamente en una
discontinuidad.}
Como $\Icl_a$ solo toma valores racionales $m/n$ con $n\le k$, la dinámica es
constante a trozos en $\theta_i$ y salta en esos racionales. En las figuras~1 y
2, $\theta=(0{,}5-0{,}4)/(0{,}6-0{,}4)=1/2$, justo en un salto: un agente con
4 vecinos iguales y 4 distintos está exactamente en el umbral. Además, en
aritmética de coma flotante la suma $0{,}4+\dots+0{,}6$ sobre 8 vecinos da
$0{,}49999999999999994$, $0{,}5$ o $0{,}5000000000000001$ según el orden en que se
recorran los vecinos (lo hemos comprobado), así que la decisión ``cómodo/incómodo''
depende del orden de iteración. Recomendaciones: (i) comparar recuentos enteros
en lugar de promedios en coma flotante; (ii) en los barridos, muestrear
$\theta$ fuera de la retícula de racionales o estudiar explícitamente las
mesetas.

\paragraph{O6. Regla de movimiento subespecificada.}
El borrador dice que el agente ``busca una casilla disponible en la que su
incomodidad sea menor''. El código implementa algo concreto: la casilla vacía
\emph{más cercana} en distancia de Chebyshev donde mejora, elegida al azar entre
las de ese radio; actualización asíncrona en orden aleatorio por barrido. Otras
elecciones razonables (vacío aleatorio global con mejora, mejor casilla global,
movimiento restringido a radio $R$) cambian la dinámica cuantitativa y a veces
la cualitativa. Hay que fijar y declarar la regla (idealmente con el protocolo
ODD \cite{Grimm2020}) y tratar la regla como un factor experimental secundario.

\paragraph{O7. Errores y puntos a corregir en la sección 6.}
\begin{itemize}
  \item La conclusión ``$h\le C_{-1,-1}+C_{1,1}$ y por ende
  $h=C_{-1,-1}+C_{1,1}$'' es incorrecta: una desigualdad no implica la igualdad.
  Lo correcto es que $h=\frac{N_1}{N}C_{1,1}+\frac{N_{-1}}{N}C_{-1,-1}$ es un
  promedio ponderado, luego
  $\min(C_{1,1},C_{-1,-1})\le h\le\max(C_{1,1},C_{-1,-1})$.
  \item La matriz $C$ se define como $3\times3$ con una fila $C_{\vac,\cdot}$,
  pero los vacíos no son agentes y esa fila no está definida. Además
  $C_{i,\vac}$ se normaliza por $n_a$ (vecinos ocupados), con lo que
  $C_{i,1}+C_{i,-1}+C_{i,\vac}\neq1$. Conviene definir dos matrices de
  contacto: $C^{\mathrm{occ}}$ (normalizada por $n_a$, natural para \modP) y
  $C^{\mathrm{full}}$ (normalizada por $k=|V|$, incluyendo la columna de
  vacíos, natural para \modS).
  \item La identidad $\bar I^H_i=\sum_j C_{i,j}H_{i,j}$ es válida para \modP{}
  con la convención de aislamiento $0$; para \modS{} la identidad análoga usa
  los recuentos medios: $\bar I^H_i=\sum_j \bar n_{i,j}H_{i,j}$.
  \item En la sección 2 se escribe $n_{a,\vac}=8-n_a$: debe ser $k-n_a$, con
  $k$ el tamaño de la vecindad, para que valga en cualquier geometría.
\end{itemize}

\paragraph{O8. Figuras.}
Ninguna figura indica tamaño de la retícula, densidad de vacíos, composición,
semilla, número de pasos, regla de movimiento ni programa utilizado; la leyenda
de colores es ilegible. Cada figura es una sola realización, por lo que no
permite afirmaciones cuantitativas. Las figuras deberían sustituirse por
(i) instantáneas acompañadas de sus parámetros completos y (ii) gráficas de
observables promediados sobre muchas realizaciones.

\paragraph{O9. Posicionamiento en la literatura.}
La introducción está vacía. Hay al menos cuatro líneas con las que el trabajo
debe dialogar explícitamente, porque delimitan qué es nuevo:
\begin{itemize}
  \item \textbf{Sakoda} \cite{Sakoda1971,Hegselmann2017} y su reanálisis
  \cite{Goles2017,Goles2020}: el modelo de Sakoda ya usa una matriz de
  actitudes (atracción, indiferencia, repulsión) entre grupos, con los vacíos
  como referencia. El modelo \modS{} es, esencialmente, un modelo de Sakoda con
  actitudes continuas y umbral de Schelling. Hay que decir con precisión qué
  añadimos.
  \item \textbf{Física estadística} \cite{Vinkovic2006,DallAsta2008,
  Gauvin2009,Grauwin2009,Castellano2009}: analogías con sistemas de espines,
  diagramas de fases en (tolerancia, densidad de vacíos), existencia o no de
  funciones potenciales. Con dos tipos y vacíos, \modS{} con interacciones
  simétricas es un gas reticular de tipo Blume--Emery--Griffiths
  \cite{BEG1971} a temperatura cero con saltos no locales.
  \item \textbf{Juegos de Schelling} \cite{Chauhan2018,Echzell2019,
  Elkind2019}: estudian la convergencia de dinámicas de mejora en versiones
  estratégicas del modelo; los resultados de potencial de la
  sección~\ref{sec:marco} son de este tipo \cite{MondererShapley1996}.
  \item \textbf{Heterogeneidad de tolerancias}
  \cite{HatnaBenenson2012,HatnaBenenson2015,AcunaGoles2025}: por O2, parte de
  lo que se obtiene en \modP{} ya está en esta línea.
\end{itemize}

\paragraph{O10. Generalidad y conexión con el trabajo previo del grupo.}
El borrador fija dos tipos $\{1,-1\}$ (el código usa $\{1,2\}$). La
formulación debería hacerse para $n$ tipos desde el principio, porque (a) el
código ya tiene la versión de tres tipos, y (b) el modelo del artículo sobre
desplazamiento de la población neutral (\texttt{modelo.py}) es un caso
particular de \modS{} con la matriz $3\times3$ (orden $1,0,-1$)
\[
  H^{\mathrm{neu}}=\begin{pmatrix}0&0&1\\1&0&1\\1&0&0\end{pmatrix},
\]
lo que da continuidad al programa y permite validar la implementación nueva
contra resultados ya publicados.

\paragraph{Menores.}
Autor ``pepito''; erratas (``llamadl'', ``qye'', ``uqe'', ``neotral'',
``Cuadnod'', ``interprestar'', ``aparecien'', ``lta'', ``modod'',
``imiytando'', ``siguinete'', ``independientee'', ``sul'', ``tota'', ``ue'');
una nota interna en el cuerpo del texto (``el programa se demora en sacar la
imagen''); unificar la terminología (``matriz de incomodidad'' en el borrador,
``matriz de aversión'' en el código, ``afinidad--aversión'' en la descripción
del proyecto).

\subsection{Observaciones sobre el código}

\paragraph{C1. Autoconteo al evaluar casillas candidatas (posible sesgo).}
En \texttt{move()}, la incomodidad en una casilla candidata $y$ se calcula con
\texttt{get\_neighbors(y)} mientras el agente sigue en su posición $x$. Si $y$
es adyacente a $x$ (siempre ocurre en la búsqueda de radio 1), el propio agente
se cuenta como vecino de sí mismo en $y$, con peso $H_{i,i}$. Con
$H^{\mathrm{cl}}$ el efecto es nulo porque $H_{i,i}=0$, pero con matrices
generales sesga la elección: si $H_{i,i}<0$, sobrevalora los destinos
adyacentes; si $H_{i,i}>0$, los infravalora. También impide que la identidad de
potencial (Proposición~\ref{prop:potencial}) se cumpla exactamente. Debe
corregirse excluyendo al agente de su propio vecindario.

\paragraph{C2. Coste de la búsqueda.}
Un agente incómodo que no puede mejorar recorre toda la retícula en cada paso,
con coste $O(L^2k)$. En los regímenes de búsqueda perpetua (O4) esto domina el
tiempo de cálculo e impide barridos amplios. Propuesta: un motor vectorizado en
\texttt{numpy} que mantenga, para cada tipo $j$, el campo de recuentos
$n_j(y)$ en todas las casillas (convolución con el núcleo de la vecindad) y lo
actualice localmente tras cada movimiento; la evaluación de un candidato pasa a
ser una consulta $O(1)$. Mesa se mantiene para la visualización.

\paragraph{C3. Tipos de estado final.}
El criterio de parada (ningún agente se movió en un barrido) no distingue entre
un equilibrio \emph{satisfecho} (todos cómodos) y un equilibrio
\emph{bloqueado} (hay incómodos que no encuentran mejora). Hay que registrar
ambos, porque en los regímenes heterofílicos y mixtos son distintos. Tampoco hay
límite de pasos ni detección de ciclos.

\paragraph{C4. Geometría fija.}
La vecindad está fijada a Moore de radio 1 (\texttt{moore=True}); además la
vecindad de interacción y la de búsqueda coinciden en forma. Deben ser
parámetros independientes: tipo (Moore/von Neumann), radio de interacción $r$ y
regla/radio de búsqueda.

\paragraph{C5. Inconsistencias menores.}
En \texttt{schelling\_general3.py} la documentación dice que la matriz por
defecto tiene unos fuera de la diagonal, pero los valores por defecto de
\texttt{\_\_init\_\_} forman un \emph{ciclo} de aversiones
($1\to2\to3\to1$), y los valores de reserva de \texttt{\_flotante\_valido}
difieren de ambos. El ciclo es, de hecho, un caso muy interesante
(Conjetura~\ref{conj:ciclo}), pero debe ser una elección explícita.

% ===============================================================
\section{Marco unificado propuesto}
\label{sec:marco}
% ===============================================================

\subsection{Definiciones}

\begin{definicion}[Espacio, tipos y vecindades]
Sea $G=(V,E)$ el grafo de sitios: el toro $\mathbb{Z}_L^2$ con la vecindad
$V_r^{\mathrm M}(x)=\{y\neq x:\|y-x\|_\infty\le r\}$ (Moore, $k=(2r+1)^2-1$) o
$V_r^{\mathrm N}(x)=\{y\neq x:\|y-x\|_1\le r\}$ (von Neumann, $k=2r(r+1)$), o
bien una red. Hay $n$ tipos $T=\{1,\dots,n\}$ y un símbolo $\vac$ para el vacío.
Una configuración es $\eta:V\to T\cup\{\vac\}$ con $N_i$ sitios de tipo $i$ y
densidad de vacíos $\rho_\vac$. Para un sitio $x$ y un agente $a$ (que se excluye
de su propio vecindario), $n_j(x)$ es el número de vecinos de $x$ de tipo $j$,
$n(x)=\sum_{j\in T}n_j(x)$ y $n_\vac(x)=k-n(x)$.
\end{definicion}

\begin{definicion}[Matriz extendida]
La matriz de afinidad--aversión extendida es
$H\in\R^{n\times(n+1)}$, con columnas indexadas por $T\cup\{\vac\}$:
$H_{i,j}$ es la incomodidad que un agente de tipo $i$ asocia a un vecino de tipo
$j$, y $H_{i,\vac}$ la que asocia a una casilla vecina vacía. Cada tipo tiene
un umbral $\tau_i$ (el caso del borrador es $\tau_i\equiv\tau$).
\end{definicion}

\begin{definicion}[Modelos \modP{} y \modS]
Para un agente de tipo $i$ situado (o evaluado) en $x$:
\begin{align*}
 \modP:&\quad I^{\modP}_i(x)=
 \begin{cases}\dfrac{1}{n(x)}\displaystyle\sum_{j\in T}n_j(x)H_{i,j},& n(x)>0,\\[2mm]
 \iota_i,& n(x)=0,\end{cases}\\
 \modS:&\quad I^{\modS}_i(x)=\sum_{j\in T}n_j(x)H_{i,j}+n_\vac(x)H_{i,\vac}.
\end{align*}
El código actual corresponde a \modS{} con $H_{i,\vac}=0$; el borrador a
\modP{} con $\iota_i=0$. A densidad máxima ($n(x)=k$) se tiene
$I^{\modS}_i=k\,I^{\modP}_i$: ambos modelos coinciden y la diferencia entre
ellos es exactamente el tratamiento de los vacíos.
\end{definicion}

\begin{definicion}[Dinámica]
Actualización asíncrona: en cada barrido los agentes se recorren en orden
aleatorio. Un agente $a$ de tipo $i$ en $x$ es \emph{cómodo} si
$I_i(x)\le\tau_i$. Si es incómodo, elige un destino $y$ en el conjunto de
vacíos que \emph{mejoran}, $\{y:\eta(y)=\vac,\ I_i(y)<I_i(x)\}$, según una
regla de búsqueda $\mathcal B$ (por defecto: los de menor distancia a $x$,
uniformemente al azar entre ellos). Si el conjunto es vacío, permanece. El
estado final es \emph{satisfecho} si todos son cómodos y \emph{bloqueado} si hay
incómodos sin mejora disponible.
\end{definicion}

\subsection{Invariancias}

\begin{proposicion}[Invariancia en \modP]\label{prop:invP}
Para cada tipo $i$, la transformación de la fila
$H_{i,\cdot}\mapsto\lambda_iH_{i,\cdot}+\mu_i$, junto con
$\tau_i\mapsto\lambda_i\tau_i+\mu_i$ e $\iota_i\mapsto\lambda_i\iota_i+\mu_i$,
con $\lambda_i>0$, deja invariante la dinámica. Si $\tau$ es común, solo las
transformaciones globales ($\lambda_i=\lambda$, $\mu_i=\mu$) son invariancias,
y solo si $\iota$ se transforma también.
\end{proposicion}

\begin{proposicion}[Invariancia en \modS]\label{prop:invS}
En un grafo $k$-regular, la transformación de la fila extendida (incluida la
columna $\vac$) $H_{i,\cdot}\mapsto\lambda_iH_{i,\cdot}+\mu_i$, con
$\tau_i\mapsto\lambda_i\tau_i+k\mu_i$ y $\lambda_i>0$, deja invariante la
dinámica. Sumar una constante solo a las columnas ocupadas \emph{no} es una
invariancia: equivale a cambiar $H_{i,\vac}$, es decir, el atractivo relativo
de los vacíos. En grafos no regulares (redes de Erdős--Rényi) ni siquiera la
traslación de la fila completa es una invariancia, porque $k$ depende del sitio.
\end{proposicion}

Ambas proposiciones son inmediatas. Su interés es que fijan los
\emph{parámetros efectivos} de cada modelo, que son los que hay que barrer.

\subsection{Reducción del modelo \modP{} con dos poblaciones}

\begin{proposicion}[Reducción]\label{prop:reduccion}
Sea $n=2$, $\bar\imath$ el otro tipo, $c_i=H_{i,\bar\imath}-H_{i,i}\neq0$,
$f(x)=n_{\bar\imath}(x)/n(x)$ la fracción de vecinos distintos y $g=1-f$ la de
iguales. Definamos
\[
 s_i=\operatorname{sgn}c_i,\qquad
 \theta_i=\begin{cases}\dfrac{\tau_i-H_{i,i}}{c_i},&c_i>0,\\[2mm]
 \dfrac{\tau_i-H_{i,\bar\imath}}{|c_i|},&c_i<0,\end{cases}\qquad
 \xi_i=\begin{cases}\dfrac{\iota_i-H_{i,i}}{c_i},&c_i>0,\\[2mm]
 \dfrac{\iota_i-H_{i,\bar\imath}}{|c_i|},&c_i<0.\end{cases}
\]
Entonces, en \modP, un agente de tipo $i$ con $s_i=+1$ se comporta como un
agente del modelo clásico con umbral $\theta_i$ sobre la fracción de distintos
$f$; con $s_i=-1$, como un agente clásico ``invertido'' con umbral $\theta_i$
sobre la fracción de iguales $g$. En ambos casos una casilla aislada equivale a
una casilla ficticia con valor $\xi_i$. En consecuencia, la dinámica de \modP{}
con $n=2$ queda determinada por $(s_i,\theta_i,\xi_i)_{i=1,2}$, la geometría,
las densidades y la regla de búsqueda.
\end{proposicion}

\begin{proof}[Esbozo]
Si $c_i>0$, $I=H_{i,i}+c_if$ es afín creciente en $f$. Si $c_i<0$,
$I=H_{i,\bar\imath}+|c_i|\,g$ es afín creciente en $g$. Se aplica la
Proposición~\ref{prop:invP} fila a fila.
\end{proof}

\begin{corolario}[Regímenes por población]
Para cada tipo: si $\theta_i\ge1$ la población es \emph{inerte} (nunca
incómoda); si $\theta_i<0$ está en \emph{búsqueda perpetua} (siempre incómoda);
si $0\le\theta_i<1$ es \emph{activa}. Si $c_i=0$ es neutral y, según
$\tau_i\gtrless H_{i,i}$, inerte o en búsqueda perpetua. Combinando
orientación y régimen, el espacio de parámetros de \modP{} con $n=2$ se reduce a
cuatro cuadrados $(\theta_1,\theta_2)\in[0,1)^2$, uno por par de orientaciones
(HH, HX, XH, XX, con H homofílica y X heterofílica), más los bordes triviales.
\end{corolario}

\begin{observacion}
La novedad genuina de la formulación matricial en \modP{} aparece con
$n\ge3$: la incomodidad
$I_i=H_{i,i}+\sum_{j\neq i}(H_{i,j}-H_{i,i})\,p_j$ depende de un vector de
contrastes de dimensión $n-1$ y ya no se reduce a una sola fracción.
\end{observacion}

\subsection{El modelo \modS{} con dos poblaciones: arquetipos de preferencia}

Con $a_i=H_{i,i}-H_{i,\vac}$ y $b_i=H_{i,\bar\imath}-H_{i,\vac}$ se tiene
$I^{\modS}_i(x)-kH_{i,\vac}=a_i\,n_i(x)+b_i\,n_{\bar\imath}(x)$: la preferencia
es una forma lineal sobre los puntos enteros del triángulo
$\{(n_i,n_{\bar\imath}):n_i+n_{\bar\imath}\le k\}$. Por la
Proposición~\ref{prop:invS}, con umbrales por tipo solo importan la dirección
$\varphi_i=\operatorname{atan2}(b_i,a_i)$ y el umbral reducido
$\hat\tau_i=(\tau_i-kH_{i,\vac})/\|(a_i,b_i)\|$. La circunferencia de
direcciones contiene arquetipos cualitativamente distintos, muchos de los cuales
no existen en \modP:

\begin{center}
\small
\begin{tabular}{@{}ccl@{}}
\toprule
$a_i$ (iguales) & $b_i$ (distintos) & arquetipo \\
\midrule
$0$ & $+$ & Schelling clásico (código actual): indiferente a iguales y vacíos \\
$-$ & $+$ & segregacionista de Sakoda: atraído por iguales, repelido por distintos \\
$-$ & $0$ & cohesivo: busca iguales, indiferente a distintos \\
$-$ & $-$ & gregario: busca densidad (con preferencia según $a_i\lessgtr b_i$) \\
$+$ & $+$ & misántropo: busca aislamiento \\
$+$ & $-$ & heterofílico: busca distintos, huye de iguales \\
$0$ & $-$ & busca distintos, indiferente a iguales \\
$+$ & $0$ & huye de iguales, indiferente a distintos \\
\bottomrule
\end{tabular}
\end{center}

Aquí la asimetría $b_1\neq b_2$ \emph{sí} tiene contenido dinámico, y el signo
relativo de $b_1$ y $b_2$ resulta decisivo.

\subsection{Funciones potenciales y convergencia en \modS}

\begin{proposicion}[Potencial ponderado]\label{prop:potencial}
Supongamos que existen pesos $w_i>0$ y una matriz simétrica $S\in\R^{n\times n}$
tales que $H_{i,j}-H_{i,\vac}=w_iS_{i,j}$ para todo $i,j\in T$. Sea
\[
 \Phi(\eta)=\sum_{\{u,v\}\in E,\ \eta(u),\eta(v)\neq\vac}S_{\eta(u),\eta(v)} .
\]
Si un agente $a$ de tipo $i$ se mueve de $x$ a $y$ (excluyéndose a sí mismo de
su vecindario), entonces
$\Phi(\eta')-\Phi(\eta)=\bigl(I_i(y)-I_i(x)\bigr)/w_i$. En consecuencia, toda
trayectoria de la dinámica \modS{} alcanza un punto fijo en un número finito de
movimientos, para cualquier umbral, regla de búsqueda y orden de
actualización.
\end{proposicion}

\begin{proof}[Esbozo]
$I_i(x)=kH_{i,\vac}+w_i\sum_{v\in V(x),\,\eta(v)\neq\vac}S_{i,\eta(v)}$.
Al mover $a$ solo cambian las aristas incidentes en $a$: se pierden las de $x$ y
se ganan las de $y$, que es exactamente $(I_i(y)-I_i(x))/w_i$. Como cada
movimiento es una mejora estricta, $\Phi$ decrece estrictamente y el espacio de
configuraciones es finito \cite{MondererShapley1996}.
\end{proof}

\begin{corolario}
Con $n=2$ la condición es $b_1b_2>0$ o $b_1=b_2=0$, sin restricción sobre
$a_1,a_2$ ni simetría de $H$. En particular, el modelo clásico de suma converge
siempre y su potencial es el número de aristas entre tipos distintos. El único
caso de dos poblaciones sin convergencia garantizada es el de
\emph{persecución}, $b_1b_2<0$: una población es atraída por la otra, que la
rehúye.
\end{corolario}

\begin{conjetura}[Potencial lexicográfico]\label{conj:lex}
Si el grafo dirigido ``$i$ se preocupa por $j$'' ($i\to j$ si
$H_{i,j}\neq H_{i,\vac}$), una vez agrupadas las clases que satisfacen la
condición de la Proposición~\ref{prop:potencial}, es acíclico, existe un
potencial lexicográfico y la dinámica converge. Ejemplos: $n=2$ con $b_1=0$; y
el modelo de población neutral $H^{\mathrm{neu}}$, donde los extremos son
indiferentes a los neutrales ($\Phi_1=$ aristas $(1,-1)$ decrece con los
movimientos de los extremos y no cambia con los de los neutrales;
$\Phi_2=$ incomodidad total de los neutrales decrece con los movimientos de
estos). Esto daría una prueba de convergencia para el modelo del artículo sobre
la población neutral.
\end{conjetura}

\begin{conjetura}[Ciclos de aversión]\label{conj:ciclo}
Con $n=3$ y aversiones cíclicas ($1$ rehúye a $2$, $2$ a $3$, $3$ a $1$, resto
indiferente), la condición de ciclo
$H_{12}H_{23}H_{31}=H_{13}H_{32}H_{21}$ necesaria para un potencial ponderado
falla. Conjeturamos que existen trayectorias que no convergen (dinámica
persistente tipo piedra--papel--tijera) en un rango abierto de umbrales y
densidades.
\end{conjetura}

\begin{observacion}
En \modP{} no hay en general función potencial, ni siquiera con $H$ simétrica,
porque la normalización por $n(x)$ rompe la simetría de las contribuciones de
cada arista. La convergencia en \modP{} es por tanto una cuestión abierta que
hay que estudiar numéricamente (tasa de no convergencia, detección de ciclos) y,
si es posible, con contraejemplos explícitos en retículos pequeños.
\end{observacion}

% ===============================================================
\section{Observables y protocolo experimental}
\label{sec:protocolo}
% ===============================================================

\subsection{Observables}

\begin{longtable}{@{}p{4.2cm}p{10.3cm}@{}}
\toprule
Observable & Definición y uso \\
\midrule
\endhead
Fracción de incómodos $u_i(t)$ & Por tipo; su valor final distingue estados
satisfechos ($u=0$) y bloqueados ($u>0$). \\
Incomodidad media $\bar I_i(t)$ & Por tipo, en la escala de cada modelo. \\
Matrices de contacto & $C^{\mathrm{occ}}_{i,j}$ (normalizada por $n(x)$) y
$C^{\mathrm{full}}_{i,j}$, $j\in T\cup\{\vac\}$ (normalizada por $k$). \\
Homofilia $h$ & $h=\sum_i\frac{N_i}{N}C^{\mathrm{occ}}_{i,i}$. \\
Densidad de interfaz $\rho_{\mathrm{int}}$ & Fracción de aristas ocupadas entre
tipos distintos; índice de segregación
$s=1-\rho_{\mathrm{int}}/\rho_{\mathrm{int}}^{\mathrm{azar}}$ ($s>0$
segregación, $s<0$ mezcla antisegregada). \\
Parámetro de orden alternado & Para regímenes heterofílicos:
$m_{\mathrm{alt}}=\bigl|\sum_x(-1)^{x_1+x_2}\sigma(x)\bigr|/N$ (tablero de
ajedrez) y pico del factor de estructura $S(\mathbf q)$ (rayas). \\
Clústeres & Número de componentes por tipo, fracción del mayor, distribución
de tamaños. \\
Longitud característica $\ell$ & Primer cero de la función de correlación o
$1/|\mathbf q^*|$ del pico de $S(\mathbf q)$; permite comparar geometrías. \\
Dinámica & Barridos hasta el punto fijo $T_c$, número total de movimientos,
fracción de corridas no convergidas antes de $T_{\max}$, ciclos detectados
(\emph{hash} de configuraciones). \\
Potencial $\Phi(t)$ & En \modS, verificación de la Proposición~\ref{prop:potencial}
(debe ser estrictamente decreciente). \\
Aislados & Fracción de agentes con $n(x)=0$ (sensibilidad a $\iota_i$ y $H_{i,\vac}$). \\
Entropía local & Ya implementada; se conserva por continuidad. \\
\bottomrule
\end{longtable}

\subsection{Protocolo de simulación}

\begin{itemize}
  \item \textbf{Tamaños:} $L\in\{40,80,160\}$ (análisis de tamaño finito en la
  fase~5); composición $N_1=N_2$ salvo que se indique; $\rho_\vac\in\{0{,}05,
  0{,}1,0{,}2,0{,}3\}$.
  \item \textbf{Réplicas:} $R=50$--$100$ semillas por punto; media e
  intervalo de confianza del 95\,\% por \emph{bootstrap}.
  \item \textbf{Parada:} punto fijo, ciclo detectado o $T_{\max}$ barridos;
  siempre se registra la causa.
  \item \textbf{Aritmética exacta:} comparaciones con recuentos enteros; los
  umbrales se muestrean evitando la retícula de racionales $m/n$, $n\le k$, o se
  estudian las mesetas explícitamente (O5).
  \item \textbf{Reproducibilidad:} configuración en ficheros YAML; cada salida
  guarda parámetros, semilla y \emph{hash} de la versión del código; resultados en
  formato tabular (Parquet o CSV) con una fila por corrida y series temporales
  submuestreadas.
\end{itemize}

% ===============================================================
\section{Programa de investigación por fases}
\label{sec:programa}
% ===============================================================

\begin{fase}{Fase 0. Especificación e implementación unificada}
\item[Objetivo.] Una sola implementación, verificada, de \modP{} y \modS{} para
$n$ tipos y geometría parametrizada.
\item[Tareas.] (i) Especificación ODD \cite{Grimm2020} del modelo unificado de
la sección~\ref{sec:marco}. (ii) Motor vectorizado con campos de recuentos
(C2), autoconteo corregido (C1), estados satisfecho/bloqueado, detección de
ciclos y $T_{\max}$ (C3), vecindades Moore/von Neumann de radio $r$ y regla de
búsqueda independiente (C4). (iii) Envoltorio Mesa para la visualización.
(iv) Pruebas automáticas: recuperación del clásico; \emph{invariancia exacta}
de trayectorias con la misma semilla bajo las transformaciones de las
Proposiciones~\ref{prop:invP}--\ref{prop:invS}; monotonía estricta de $\Phi$ en
\modS; coincidencia de \modP{} y \modS{} a densidad máxima.
\item[Entregable.] Paquete \texttt{schelling\_matricial} con pruebas y un
cuaderno de validación.
\item[Cierre.] Todas las pruebas pasan; reproducimos cualitativamente el
diagrama de fases clásico \cite{Gauvin2009} y los resultados del artículo de
la población neutral con $H^{\mathrm{neu}}$.
\end{fase}

\begin{fase}{Fase 1. Teoría del modelo \modP}
\item[Objetivo.] Reescribir las secciones 2--5 del borrador en torno a la
Proposición~\ref{prop:reduccion}.
\item[Tareas.] Demostraciones completas; tabla de regímenes
$(s_i,\text{inerte/activo/perpetuo})$; decisión sobre $\iota_i$ (O3); corrección
de la sección 6 (O7); análisis de las mesetas en $\theta_i$ (O5); búsqueda de
contraejemplos de convergencia en retículos pequeños.
\item[Entregable.] Nueva versión de las secciones teóricas del borrador.
\item[Cierre.] Revisión conjunta de los enunciados.
\end{fase}

\begin{fase}{Fase 2. Teoría del modelo \modS}
\item[Objetivo.] Desarrollar la teoría paralela para la versión de suma.
\item[Tareas.] Proposición~\ref{prop:invS}; clasificación por direcciones
$\varphi_i$; prueba de la Proposición~\ref{prop:potencial} y de la
Conjetura~\ref{conj:lex}; cotas del número de movimientos; relación precisa con
Sakoda \cite{Goles2017} y con el modelo de Blume--Emery--Griffiths
\cite{BEG1971}.
\item[Entregable.] Sección teórica nueva ``El modelo de suma''.
\item[Cierre.] Pruebas escritas y verificadas numéricamente (fase 0).
\end{fase}

\begin{fase}{Fase 3. Diagramas de régimen con dos poblaciones (Moore, $r=1$)}
\item[Objetivo.] Mapas cuantitativos de la dinámica en la geometría de
referencia.
\item[Tareas.] \modP: barrido de $(\theta_1,\theta_2)\in[0,1)^2$ para los
cuatro pares de orientaciones y cuatro valores de $\rho_\vac$. \modS: barrido de
$(\varphi_1,\varphi_2)$ en la circunferencia con $\hat\tau_i$ fijo, y de
$\hat\tau$ en direcciones representativas; atención especial al régimen de
persecución $b_1b_2<0$.
\item[Entregable.] Mapas de $s$, $u$, $T_c$, $\ell$ y tasa de no convergencia;
localización de transiciones.
\item[Cierre.] Mapas estables al duplicar $R$ y al pasar de $L=40$ a $L=80$.
\end{fase}

\begin{fase}{Fase 4. Comparación \modP{} frente a \modS}
\item[Objetivo.] Cuantificar el efecto de normalizar por el número de vecinos.
\item[Tareas.] Emparejar parámetros equivalentes a densidad máxima y variar
$\rho_\vac$; estudiar el papel de $H_{i,\vac}$ e $\iota_i$.
\item[Entregable.] Figura de divergencia entre modelos en función de
$\rho_\vac$.
\item[Cierre.] Contraste de la Hipótesis~\ref{hip:densidad}.
\end{fase}

\begin{fase}{Fase 5. Geometría}
\item[Objetivo.] Efecto del tipo y radio de la vecindad.
\item[Tareas.] Moore y von Neumann con $r\in\{1,2,3\}$ ($k=8,24,48$ y
$k=4,12,24$); umbrales de \modS{} escalados con $k$; análisis de tamaño finito
en $L$; regímenes heterofílicos en vecindades bipartitas y no bipartitas.
\item[Entregable.] Leyes de escala de $\ell$ con $r$; comparación de patrones
alternados.
\item[Cierre.] Contraste de las Hipótesis~\ref{hip:frustracion} y
\ref{hip:escala}.
\end{fase}

\begin{fase}{Fase 6. Tres poblaciones}
\item[Objetivo.] Explorar lo que no se reduce al caso clásico.
\item[Tareas.] Matrices jerárquicas (generalizando $H^{\mathrm{neu}}$) frente a
cíclicas; prueba o refutación de la Conjetura~\ref{conj:ciclo}; espacio de
contrastes de dimensión 2 en \modP.
\item[Entregable.] Catálogo de regímenes para $n=3$.
\item[Cierre.] Contraste de las Hipótesis~\ref{hip:jerarquia} y
\ref{hip:ciclo}.
\end{fase}

\begin{fase}{Fase 7 (opcional). Redes y heterogeneidad}
\item[Objetivo.] Extensiones naturales ya parcialmente implementadas.
\item[Tareas.] Redes de Erdős--Rényi (\texttt{schelling\_red2.py}), recordando que
en \modS{} la traslación deja de ser invariancia; umbrales individuales
\cite{AcunaGoles2025}.
\item[Entregable.] Nota o sección adicional.
\end{fase}

\begin{fase}{Fase 8. Redacción}
\item[Objetivo.] Convertir el trabajo en uno o dos artículos.
\item[Opciones.] (a) Un artículo teórico--numérico con \modP{} y \modS{} para
$n=2$ y la geometría; (b) un segundo artículo sobre $n=3$ y convergencia. La
decisión se toma al cerrar la fase 4.
\end{fase}

% ===============================================================
\section{Hipótesis de trabajo}
% ===============================================================

\begin{hipotesis}[Equivalencia en \modP]
En \modP{} con $n=2$, dos matrices con los mismos $(s_i,\theta_i,\xi_i)$
producen distribuciones idénticas de todos los observables; en particular la
asimetría $A=\beta-\delta$ no tiene efecto propio.
\end{hipotesis}

\begin{hipotesis}[Coincidencia a alta densidad]\label{hip:densidad}
La discrepancia entre \modP{} y \modS{} con parámetros emparejados crece con
$\rho_\vac$ y se anula cuando $\rho_\vac\to0$.
\end{hipotesis}

\begin{hipotesis}[Persecución]
En \modS{} con $b_1b_2<0$ existe una región de umbrales en la que la fracción
de corridas no convergidas no tiende a cero al aumentar $T_{\max}$.
\end{hipotesis}

\begin{hipotesis}[Frustración geométrica]\label{hip:frustracion}
En regímenes heterofílicos, la vecindad de von Neumann de radio 1 (grafo
bipartito) produce orden de tablero de ajedrez ($m_{\mathrm{alt}}$ grande),
mientras que Moore de radio 1 y von Neumann de radio $\ge2$ (con ciclos
impares) producen rayas o laberintos con muchos estados finales degenerados y
convergencia más lenta. La figura~2 del borrador es compatible con esto.
\end{hipotesis}

\begin{hipotesis}[Escala con el radio]\label{hip:escala}
En regímenes homofílicos, la longitud característica de los dominios crece
linealmente con $r$, y el umbral de aparición de segregación, expresado en
$\theta$, depende débilmente de $k$ salvo por los saltos en racionales.
\end{hipotesis}

\begin{hipotesis}[Jerarquía]\label{hip:jerarquia}
Las matrices de $n=3$ con grafo de preocupación acíclico convergen siempre en
\modS.
\end{hipotesis}

\begin{hipotesis}[Ciclos]\label{hip:ciclo}
Las matrices cíclicas de $n=3$ presentan dinámica persistente en una región
abierta de parámetros.
\end{hipotesis}

% ===============================================================
\section{Calendario orientativo y dependencias}
% ===============================================================

\begin{center}
\begin{tabular}{@{}llll@{}}
\toprule
Fase & Depende de & Duración estimada & Producto principal \\
\midrule
0 Implementación & --- & 2--3 semanas & paquete verificado \\
1 Teoría \modP & --- (en paralelo con 0) & 2 semanas & secciones reescritas \\
2 Teoría \modS & 0 (verificación) & 3 semanas & sección nueva \\
3 Diagramas $n=2$ & 0, 1, 2 & 3--4 semanas & mapas de régimen \\
4 \modP{} frente a \modS & 3 & 2 semanas & figura de divergencia \\
5 Geometría & 3 & 3 semanas & leyes de escala \\
6 Tres poblaciones & 2, 3 & 3--4 semanas & catálogo $n=3$ \\
7 Redes (opcional) & 3 & 2 semanas & nota \\
8 Redacción & 4 (decisión), 5, 6 & 3 semanas & artículo(s) \\
\bottomrule
\end{tabular}
\end{center}

\section{Decisiones pendientes para el equipo}

\begin{enumerate}
  \item Convención de aislamiento en \modP: ¿$\iota_i$ libre, $\iota_i=H_{i,i}$
  (el aislamiento vale como estar entre iguales) u otra?
  \item Columna de vacíos en \modS: ¿mantener $H_{i,\vac}=0$ como en el código o
  tratarla como parámetro?
  \item ¿Umbral común $\tau$ o umbrales por tipo $\tau_i$? (Los segundos hacen
  exactas las invariancias fila a fila y no complican la implementación.)
  \item Regla de búsqueda por defecto (la actual: vacío mejorante más cercano).
  \item Alcance de la publicación (fase 8).
\end{enumerate}

\section{Primer paso}

Proponemos empezar por la fase 0 y, en paralelo, la fase 1: (i) escribir la
especificación ODD del modelo unificado y acordar las decisiones de la sección
anterior; (ii) corregir el autoconteo (C1) en el código actual, porque afecta a
cualquier resultado con $H_{i,i}\neq0$; (iii) implementar el motor vectorizado
con sus pruebas de invariancia y de potencial. Con eso tendremos una base
fiable para las fases numéricas.

% ===============================================================
\begin{thebibliography}{99}
% ===============================================================

\bibitem{Schelling1971} T.~C. Schelling. Dynamic models of segregation.
\emph{Journal of Mathematical Sociology}, 1(2):143--186, 1971.

\bibitem{Schelling1978} T.~C. Schelling. \emph{Micromotives and
Macrobehavior}. W.~W. Norton, Nueva York, 1978.

\bibitem{Sakoda1971} J.~M. Sakoda. The checkerboard model of social
interaction. \emph{Journal of Mathematical Sociology}, 1(1):119--132, 1971.

\bibitem{Hegselmann2017} R.~Hegselmann. Thomas C. Schelling and James M.
Sakoda: The intellectual, technical, and social history of a model.
\emph{JASSS}, 20(3):15, 2017.

\bibitem{Goles2017} E.~Goles, N.~G. Domic, S.~Rica. Self-organized
societies: On the Sakoda model of social interactions. \emph{Complexity},
2017:3548591, 2017.

\bibitem{Goles2020} E.~Goles, S.~Rica, G.~A. Ruz. From immigration to social
segregation: A view through the Sakoda model of social interactions.
\emph{Scientific Reports}, 10, 2020.

\bibitem{HatnaBenenson2012} E.~Hatna, I.~Benenson. The Schelling model of
ethnic residential dynamics: Beyond the integrated--segregated dichotomy of
patterns. \emph{JASSS}, 15(1):6, 2012.

\bibitem{HatnaBenenson2015} E.~Hatna, I.~Benenson. Combining segregation and
integration: Schelling model dynamics for heterogeneous population.
\emph{JASSS}, 18(4):15, 2015.

\bibitem{AcunaGoles2025} G.~Acuña, E.~Goles. Individual tolerances and
dynamics of intragroup segregation in an extended Schelling model.
\emph{Scientific Reports}, 15:43050, 2025.

\bibitem{Castellano2009} C.~Castellano, S.~Fortunato, V.~Loreto. Statistical
physics of social dynamics. \emph{Reviews of Modern Physics}, 81(2):591--646,
2009.

\bibitem{Vinkovic2006} D.~Vinković, A.~Kirman. A physical analogue of the
Schelling model. \emph{PNAS}, 103(51):19261--19265, 2006.

\bibitem{DallAsta2008} L.~Dall'Asta, C.~Castellano, M.~Marsili. Statistical
physics of the Schelling model of segregation. \emph{Journal of Statistical
Mechanics}, L07002, 2008.

\bibitem{Gauvin2009} L.~Gauvin, J.~Vannimenus, J.-P. Nadal. Phase diagram of
a Schelling segregation model. \emph{European Physical Journal B},
70:293--304, 2009.

\bibitem{Grauwin2009} S.~Grauwin, E.~Bertin, R.~Lemoy, P.~Jensen. Competition
between collective and individual dynamics. \emph{PNAS},
106(49):20622--20626, 2009.

\bibitem{BEG1971} M.~Blume, V.~J. Emery, R.~B. Griffiths. Ising model for the
$\lambda$ transition and phase separation in He$^3$--He$^4$ mixtures.
\emph{Physical Review A}, 4(3):1071--1077, 1971.

\bibitem{MondererShapley1996} D.~Monderer, L.~S. Shapley. Potential games.
\emph{Games and Economic Behavior}, 14(1):124--143, 1996.

\bibitem{Chauhan2018} A.~Chauhan, P.~Lenzner, L.~Molitor. Schelling
segregation with strategic agents. En \emph{SAGT 2018}, LNCS 11059,
pp.~137--149. Springer, 2018.

\bibitem{Echzell2019} H.~Echzell, T.~Friedrich, P.~Lenzner, L.~Molitor,
M.~Pappik, F.~Schöne, F.~Sommer, D.~Stangl. Convergence and hardness of
strategic Schelling segregation. En \emph{WINE 2019}, LNCS 11920,
pp.~156--170. Springer, 2019.

\bibitem{Elkind2019} E.~Elkind, J.~Gan, A.~Igarashi, W.~Suksompong,
A.~A. Voudouris. Schelling games on graphs. En \emph{IJCAI 2019},
pp.~266--272, 2019.

\bibitem{Grimm2020} V.~Grimm et al. The ODD protocol for describing
agent-based and other simulation models: A second update to improve clarity,
replication, and structural realism. \emph{JASSS}, 23(2):7, 2020.

\bibitem{Fossett2006} M.~Fossett. Ethnic preferences, social distance
dynamics, and residential segregation. \emph{Journal of Mathematical
Sociology}, 30(3--4):185--273, 2006.

\bibitem{Zhang2011} J.~Zhang. Tipping and residential segregation: A unified
Schelling model. \emph{Journal of Regional Science}, 51(1):167--193, 2011.

\end{thebibliography}

\end{document}
